\section{Scalar tail estimates}
\label{sec:scalar}
We first record the coordinate estimates used by the comparison construction. All constants in this section may depend on the fixed tail exponents.

\begin{lemma}[Conditioning a coordinate to be small]\label{lem:small-coordinate}
Let $W=F^{-1/\alpha}T$, where $F\ge1$, $0<\alpha<2$, and $T$ has density $p_\alpha$. For $0<u\le1$,
\begin{align}
 -\log\Pp(|W|\le u)&\le C_\alpha u^{-\alpha}/F,\qquad
 \E[W\mid |W|\le u]=0,\\
 \E[W^2\mid |W|\le u]&\le C_\alpha u^{2-\alpha}/F.
 \label{eq:small-coordinate}
\end{align}
If $\alpha<1$, also
\begin{equation}
 \E[|W|\mid |W|\le u]\le C_\alpha u^{1-\alpha}/F.
 \label{eq:small-first-moment}
\end{equation}
For $|w|\le B+5$, the unconditioned density obeys
\begin{equation}
 p_{F,\alpha}(w)\ge\frac{c_\alpha'}{F(B+6)^{1+\alpha}}.
 \label{eq:density-lower}
\end{equation}
\end{lemma}
\begin{proof}
Put $v=F^{1/\alpha}u$. For $v\le1$, the small-ball probability is at least $c_\alpha'v$, and $v^\alpha\log(1/(c_\alpha'v))$ is bounded. This proves the logarithmic estimate after using $F=v^\alpha u^{-\alpha}$. Conditional absolute moments of orders one and two are at most $u$ and $u^2$; multiplying by $F$ and using $Fu^\alpha\le1$ gives the stated estimates in this range.

For $v\ge1$, the small-ball probability is bounded below, the tail is at most $C_\alpha v^{-\alpha}$, and integration gives truncated moments bounded by $C_\alpha v^{r-\alpha}$ for $r>\alpha$. These facts imply all the estimates after rescaling. Symmetry gives the conditional mean. Finally
\[
 p_{F,\alpha}(w)=\frac{c_\alpha}{F}
       (F^{-2/\alpha}+w^2)^{-(1+\alpha)/2},
\]
which proves~\eqref{eq:density-lower}.
\end{proof}

\begin{lemma}[Fractional moments of absolute row sums]\label{lem:absolute-sum}
Suppose $0<p<\alpha<1$. For independent fan-in-$F$ coordinates $W_j$ with tail index $\alpha$ and deterministic $c_j\ge0$,
\begin{equation}
 \E\left(\sum_j c_j|W_j|\right)^p
 \le C_{\alpha,p}\left(\frac1F\sum_j c_j^\alpha\right)^{p/\alpha}.
 \label{eq:absolute-sum}
\end{equation}
\end{lemma}
\begin{proof}
Write $K=F^{-1}\sum_jc_j^\alpha$. Each summand has tail at most $C_\alpha c_j^\alpha/(Fu^\alpha)$ and an unconditional truncated first moment at level $u$ at most $C_\alpha c_j^\alpha u^{1-\alpha}/F$. Split the sum according to whether any summand exceeds $u/2$. A union bound controls that event by $C_\alpha Ku^{-\alpha}$. On its complement, Markov's inequality applied to the truncated sum gives the same bound for the event that the sum exceeds $u$. Thus
\[
 \Pp\left\{\sum_jc_j|W_j|>u\right\}
 \le\min\{1,C_\alpha Ku^{-\alpha}\}.
\]
Integrate $pu^{p-1}$ times this bound to obtain~\eqref{eq:absolute-sum}. The case $K=0$ is immediate.
\end{proof}

\begin{lemma}[Unbiased distinguished coefficients]\label{lem:distinguished}
For a teacher coefficient $c\in[-B,B]$ and $0<\epsilon\le1$, there is an absolutely continuous law $Q_c$ with
\[
 \E W=c,\quad \operatorname{Var}(W)\le17\epsilon^2,\quad
 |W-c|\le9\epsilon,\quad |W|\le B+5,\quad |W|>\epsilon
 \quad\text{almost surely}.
\]
Its density is bounded by $1/(2\epsilon)$, and hence
\[
 \KL(Q_c\|p_{F,\alpha})
 \le\log\!\frac{C_\alpha F(B+6)^{1+\alpha}}\epsilon.
\]
\end{lemma}
\begin{proof}
For $|c|>4\epsilon$, use the uniform law on $[c-\epsilon,c+\epsilon]$. Otherwise mix the uniform laws on $[3\epsilon,5\epsilon]$ and $[-5\epsilon,-3\epsilon]$ with probabilities $(1+c/(4\epsilon))/2$ and $(1-c/(4\epsilon))/2$. The latter has variance at most $(16+1/3)\epsilon^2$. The support and density bounds are immediate, and the KL estimate follows from~\eqref{eq:density-lower}.
\end{proof}

The arguments use only symmetric densities comparable above and below to $(1+t^2)^{-(1+\alpha)/2}$. The same conclusions therefore hold for any such coordinate family, with modified constants. This includes symmetric $\alpha$-stable densities, whose positivity and power-law tails give the comparison; see~\cite{zolotarev}.
